% ── §12 Results: Contact Map Compressibility ──
\section{Results: Contact Maps as Boolean Functions}
\label{sec:results-compressibility}

\subsection{Motivation}
If native contact maps possess Karnaugh-map-expressible regularity beyond what
density and separation structure alone can explain, then exact
Quine--McCluskey minimisation of real contact matrices should yield fewer
prime implicants than minimisation of matched random controls.

\subsection{Method}
For six PSICOV proteins with $L \leq 62$ (the scale at which exact QM on a
padded $64 \times 64$ grid remains tractable), we minimise the ternary truth
table defined by the native contact matrix (ON = contact, OFF = candidate
non-contact, DC = all other cells including near-diagonal and padding border).
Controls shuffle contact labels within separation bins, preserving both density
and separation structure.

\subsection{Result}
All six proteins show lower compression ratios (#prime implicants / #contact
cells) for real maps than for their shuffled counterparts.  The mean
compression ratio is 2.31 for real maps versus 2.74 for shuffled controls,
with a Wilcoxon signed-rank test $p$-value of 0.031 ($6/6$ consistent
direction).  This demonstrates that protein geometry carries K-map-expressible
structure that is not explained by contact density or separation distribution
alone.
